\documentclass[a4paper,12pt]{article}

\usepackage[utf8]{inputenc}
\usepackage[T1]{fontenc}
\usepackage{lmodern}
\usepackage{amsmath}
\usepackage{amssymb}
\usepackage{bm}
\usepackage[margin=2.5cm]{geometry}

\title{\textbf{Two Rational Models for a Transient Spherical Signal}}
\author{}
\date{}

\begin{document}

\maketitle

\begin{abstract}
This note merges two closely related descriptions of a transient spherical
signal in three-dimensional space. The first case is a scalar shell model: the
signal occupies a time-dependent spherical layer whose outer radius expands
asymptotically and whose thickness eventually collapses. The second case is a
radial vector field: the signal is represented by a smooth outward field whose
magnitude is concentrated near an expanding radius. Both constructions use
rational algebraic functions rather than exponentials or Gaussian profiles.
\end{abstract}

\section{Common Setting}

Let
\[
    \bm{r}=(x,y,z)\in\mathbb{R}^3,
    \qquad
    r=\|\bm{r}\|=\sqrt{x^2+y^2+z^2},
    \qquad
    t\ge 0.
\]
When \(r>0\), the outward radial unit vector is
\[
    \hat{\bm{u}}_r=\frac{\bm{r}}{r}.
\]
At \(r=0\), this unit vector is not defined, so any vector-field model below
must either be interpreted away from the origin or completed by a convention
such as \(\bm{V}(\bm{0},t)=\bm{0}\).

The common physical picture is a signal emitted from the origin, expanding
radially, and then dissolving. The two cases differ in how the signal is
represented:
\begin{itemize}
    \item Case I describes the occupied region of space as a scalar shell.
    \item Case II describes the signal as a smooth radial vector field.
\end{itemize}

\section{Case I: Scalar Finite-Support Shell}

In the scalar shell model the signal is present in a spherical layer
\[
    R_{\mathrm{int}}(t)\le r\le R_{\mathrm{ext}}(t),
\]
and absent outside it. The outer radius expands from zero toward a finite
limit \(R_{\max}\):
\begin{equation}
    R_{\mathrm{ext}}(t)
    =
    R_{\max}\frac{t}{t+k_1},
    \qquad
    R_{\max}>0,\quad k_1>0.
    \label{eq:outer-radius}
\end{equation}
This is a rational saturation law. It begins at \(0\), increases
monotonically, and satisfies
\[
    \lim_{t\to\infty}R_{\mathrm{ext}}(t)=R_{\max}.
\]

\subsection{Thickness and Internal Radius}

The shell thickness is required to vanish both at the initial instant and at
late time. A convenient rational choice is
\begin{equation}
    h(t)
    =
    R_{\mathrm{ext}}(t)\frac{k_2}{t+k_2}
    =
    R_{\max}\frac{t}{t+k_1}\frac{k_2}{t+k_2},
    \qquad k_2>0.
    \label{eq:thickness}
\end{equation}
The internal radius is then
\begin{equation}
    R_{\mathrm{int}}(t)
    =
    R_{\mathrm{ext}}(t)-h(t)
    =
    R_{\mathrm{ext}}(t)
    \left(1-\frac{k_2}{t+k_2}\right).
    \label{eq:internal-radius}
\end{equation}
Thus \(R_{\mathrm{int}}(0)=0\), and for large \(t\) the internal radius
approaches the outer radius.

\subsection{Field Definition}

The occupied signal region is
\[
    \Omega(t)
    =
    \left\{
        \bm{r}\in\mathbb{R}^3:
        R_{\mathrm{int}}(t)\le \|\bm{r}\|\le R_{\mathrm{ext}}(t)
    \right\}.
\]
A binary scalar signal can be written with an indicator function:
\begin{equation}
    S(\bm{r},t)
    =
    \mathbf{1}_{\Omega(t)}(\bm{r}).
    \label{eq:indicator-signal}
\end{equation}
Equivalently, the support is described by
\begin{equation}
    \left(R_{\mathrm{ext}}(t)-h(t)\right)^2
    \le
    x^2+y^2+z^2
    \le
    R_{\mathrm{ext}}(t)^2.
    \label{eq:shell-support}
\end{equation}
Substituting \eqref{eq:outer-radius} and \eqref{eq:thickness}, this becomes
\begin{equation}
    \left[
        R_{\max}\frac{t}{t+k_1}
        \left(1-\frac{k_2}{t+k_2}\right)
    \right]^2
    \le
    \|\bm{r}\|^2
    \le
    \left[
        R_{\max}\frac{t}{t+k_1}
    \right]^2.
    \label{eq:explicit-shell-support}
\end{equation}

\subsection{Behavior}

For early time, \(t\ll k_1,k_2\),
\[
    R_{\mathrm{ext}}(t)
    \approx
    \frac{R_{\max}}{k_1}t,
    \qquad
    h(t)\approx R_{\mathrm{ext}}(t).
\]
The internal radius is therefore close to zero, so the signal initially behaves
like a solid sphere expanding from the origin.

For late time,
\[
    R_{\mathrm{ext}}(t)\to R_{\max},
    \qquad
    h(t)\approx R_{\max}\frac{k_2}{t}\to 0.
\]
The occupied region becomes thinner and thinner near the limiting sphere
\(r=R_{\max}\).

\section{Case II: Smooth Rational Radial Vector Field}

The second model represents the same qualitative phenomenon as a vector field
rather than as a sharply supported scalar region. The field is purely radial:
\begin{equation}
    \bm{V}(\bm{r},t)
    =
    M(r,t)\hat{\bm{u}}_r,
    \qquad r>0.
    \label{eq:vector-field-general}
\end{equation}
Here \(M(r,t)\) is a scalar magnitude concentrated near a moving shell radius.

\subsection{Rational Shell Magnitude}

The shell profile is modeled with a Lorentzian--Cauchy rational bell:
\begin{equation}
    M(r,t)
    =
    \frac{A(t)}
    {1+\left(\dfrac{r-\mu(t)}{w(t)}\right)^2}.
    \label{eq:lorentzian-magnitude}
\end{equation}
The parameters have direct interpretations:
\begin{itemize}
    \item \(\mu(t)\) is the center radius of the shell.
    \item \(w(t)\) is the shell width.
    \item \(A(t)\) is the peak amplitude.
\end{itemize}

The center radius follows another rational saturation law:
\begin{equation}
    \mu(t)
    =
    R_{\max}\frac{t}{t+k_{\mathrm{exp}}},
    \qquad k_{\mathrm{exp}}>0.
    \label{eq:center-radius}
\end{equation}
The amplitude decays rationally:
\begin{equation}
    A(t)
    =
    \frac{I_0\tau}{t+\tau},
    \qquad I_0>0,\quad \tau>0.
    \label{eq:amplitude-decay}
\end{equation}
At \(t=0\), the amplitude is \(I_0\), while for large \(t\) it decays like
\(1/t\). A simple diffusive width law is
\begin{equation}
    w(t)=w_0+\beta t,
    \qquad w_0>0,\quad \beta\ge 0.
    \label{eq:width-law}
\end{equation}

\subsection{Full Equation}

Combining \eqref{eq:vector-field-general}--\eqref{eq:width-law}, the vector
field is
\begin{equation}
    \bm{V}(\bm{r},t)
    =
    \left[
        \frac{
            \dfrac{I_0\tau}{t+\tau}
        }
        {
            1+
            \left(
                \dfrac{
                    r-R_{\max}\dfrac{t}{t+k_{\mathrm{exp}}}
                }
                {w_0+\beta t}
            \right)^2
        }
    \right]
    \frac{\bm{r}}{r},
    \qquad r>0.
    \label{eq:full-vector-field}
\end{equation}
This is not compactly supported: the Lorentzian profile has heavy tails.
However, its largest values are concentrated near
\[
    r=\mu(t),
\]
and the peak magnitude is \(A(t)\).

\subsection{Behavior}

The direction is always radial and outward when \(M(r,t)>0\). The denominator
in \eqref{eq:lorentzian-magnitude} is smallest when \(r=\mu(t)\), so the field
has a smooth shell-like maximum at the moving radius. The rational decay in
\(A(t)\) makes the field vanish in amplitude over long times.

The model is algebraically cheap: apart from the norm \(r=\|\bm{r}\|\), it uses
only arithmetic operations. In applications where avoiding a square root is
important, one may formulate approximate or squared-radius variants, but the
radial vector direction itself naturally requires normalization.

\section{Comparison of the Two Cases}

The scalar shell and vector-field models encode related but distinct ideas.

\begin{center}
\begin{tabular}{lll}
\hline
Feature & Case I & Case II \\
\hline
Quantity & scalar \(S(\bm{r},t)\) & vector \(\bm{V}(\bm{r},t)\) \\
Support & finite shell & infinite heavy-tailed profile \\
Boundary & sharp inequalities & smooth Lorentzian decay \\
Radius law & \(R_{\mathrm{ext}}(t)\) & \(\mu(t)\) \\
Dissolution & thickness \(h(t)\to 0\) & amplitude \(A(t)\to 0\) \\
Direction & none & outward radial direction \\
\hline
\end{tabular}
\end{center}

Case I is appropriate when the main object is the occupied spatial region of
the signal. Case II is appropriate when the signal should act as a force,
velocity, pressure, or displacement field.

\section{Conclusion}

Both models describe transient spherical signals using rational functions. The
first gives a sharply bounded shell whose thickness grows and then collapses.
The second gives a smooth outward radial field whose magnitude forms an
expanding rational shell and fades over time. Together they provide two
complementary descriptions of the same qualitative phenomenon: a signal born at
the origin, expanding through space, and ultimately dissolving.

\end{document}
